\documentclass[11pt]{article}

\usepackage[margin=1in]{geometry}
\usepackage{amsmath, amssymb, amsthm}
\usepackage{algorithm}
\usepackage{algpseudocode}
\usepackage{hyperref}
\usepackage{enumitem}
\usepackage{booktabs}

\title{Notes on Continuous Observation Spaces}
\author{}
\date{}

\begin{document}
\maketitle

\section{Setting}

Everything so far — Sarsa, $n$-step Sarsa, Q-learning, Monte Carlo control —
assumes a \emph{tabular} representation: $\hat{Q}(s,a)$ is looked up in a
table indexed by a finite (and usually small) set of states $s \in
\mathcal{S}$. Many environments of interest instead have a
\textbf{continuous observation space} $\mathcal{S} \subseteq \mathbb{R}^d$:
e.g.\ \texttt{MountainCar-v0} returns $(\text{position}, \text{velocity}) \in
\mathbb{R}^2$, and \texttt{CartPole} returns a 4-vector of cart position,
cart velocity, pole angle, and pole angular velocity. A table over
$\mathbb{R}^d$ is not just large — it is uncountably infinite, so no
tabular method can be applied directly, and two states that are numerically
close (and should behave almost identically) would otherwise be treated as
completely unrelated table entries.

The fix is to convert the continuous observation into a small set of
discrete features and keep using ordinary tabular control (Sarsa, in the
examples below) on top of that discretization. We cover two such feature
maps: \textbf{state aggregation} (binning) and \textbf{tile coding}
(overlapping, offset binnings).

\section{State aggregation}

State aggregation discretizes each dimension of the observation
independently into $b_i$ bins, then treats the resulting bin-index tuple
as a discrete state. For $d$-dimensional observations with bins
$b_1, \dots, b_d$, this partitions $\mathcal{S}$ into a grid of
$\prod_i b_i$ cells.

\subsection{Construction}

Given per-dimension bounds $\text{low}, \text{high} \in \mathbb{R}^d$
(available from \texttt{env.observation\_space}) and a bin count vector
$\mathbf{b} = (b_1, \dots, b_d)$, define $d$ sets of bin edges

\[
    B_i = \text{linspace}(\text{low}_i, \text{high}_i, b_i - 1), \qquad i = 1, \dots, d,
\]

i.e.\ $b_i - 1$ interior edges producing $b_i$ bins along dimension $i$.
A raw observation $o = (o_1, \dots, o_d)$ is mapped to the discrete state

\[
    s = \big(\, \mathrm{digitize}(o_1, B_1),\ \dots,\ \mathrm{digitize}(o_d, B_d) \,\big) \in \{0, \dots, b_1-1\} \times \cdots \times \{0, \dots, b_d-1\},
\]

where $\mathrm{digitize}(x, B)$ returns the index of the bin containing
$x$. This is implemented as an \texttt{ObservationWrapper}: it leaves the
environment's dynamics untouched and only rewrites the observation, so any
tabular algorithm can be pointed at it unmodified, with
$\hat{Q}$ now a table of shape $b_1 \times \cdots \times b_d \times
|\mathcal{A}|$ (e.g.\ $20 \times 20 \times 3$ for \texttt{MountainCar-v0}
with $\mathbf{b} = (20, 20)$).

\subsection{Sarsa on the aggregated environment}

Control proceeds exactly as ordinary tabular Sarsa (see the Sarsa notes),
just indexing $\hat{Q}$ by the bin tuple instead of a raw state:

\begin{algorithm}[h]
\caption{Sarsa with state aggregation}
\begin{algorithmic}[1]
\State Choose bins $\mathbf{b}$; build the aggregation wrapper from $\text{low}, \text{high}, \mathbf{b}$
\State Initialize $\hat{Q}(s,a) \leftarrow 0$ for all bin-tuples $s$ and actions $a$
\For{episode $= 1, 2, \dots$}
    \State $S_0 \leftarrow$ aggregated reset observation; choose $A_0$ $\varepsilon$-greedily from $\hat{Q}(S_0, \cdot)$
    \For{$t = 0, 1, 2, \dots$ until done}
        \State Take $A_t$ in the wrapped env, observe $R_{t+1}$ and aggregated $S_{t+1}$
        \State Choose $A_{t+1}$ $\varepsilon$-greedily from $\hat{Q}(S_{t+1}, \cdot)$
        \State $\hat{Q}(S_t, A_t) \leftarrow \hat{Q}(S_t, A_t) + \alpha \big( R_{t+1} + \gamma \hat{Q}(S_{t+1}, A_{t+1}) - \hat{Q}(S_t, A_t) \big)$
    \EndFor
\EndFor
\end{algorithmic}
\end{algorithm}

\subsection{Limitations}

Aggregation introduces a hard discontinuity at every bin edge: two
observations $o, o'$ on opposite sides of an edge but arbitrarily close in
$\mathbb{R}^d$ can map to completely different, independently-learned
table entries, while two observations far apart within the same bin are
treated as identical. Resolution is also coupled across the whole space —
finer bins improve accuracy everywhere but blow up the table size
($\prod_i b_i$ grows exponentially in $d$) and slow down learning (each
cell must be visited many times before $\hat{Q}$ is accurate there).

\section{Tile coding}

Tile coding keeps the idea of binning but overlays $n$ independent
aggregations (\textbf{tilings}) of the same space, each shifted by a
different random offset. A state is then represented not by one bin index
but by $n$ bin indices, one per tiling, and $Q(s,a)$ is estimated as the
\emph{average} of the $n$ corresponding table entries. This is a coarse
form of linear function approximation with binary features: the effective
resolution is much finer than any single tiling's bin width, because two
nearby points that fall in the same bin for most tilings but different
bins for a few tilings get almost — but not exactly — the same estimate.

\subsection{Construction}

For $n$ tilings and base bin counts $\mathbf{b}$, tiling $i = 1, \dots, n$
uses independently jittered bounds

\[
    \text{low}^{(i)} = \text{low} - u_i \cdot 0.2 \cdot \text{low}, \qquad
    \text{high}^{(i)} = \text{high} + u_i \cdot 0.2 \cdot \text{high}, \qquad u_i \sim \mathrm{Uniform}(0,1),
\]

together with a per-tiling displacement (an asymmetric offset following the
displacement vector $(1, 3, 5, \dots)$ standard in tile coding, so that no
two tilings' bin edges coincide) before computing per-dimension edges
$B_1^{(i)}, \dots, B_d^{(i)}$ exactly as in state aggregation. A raw
observation $o$ is then mapped to a list of $n$ bin-tuples,

\[
    s = \big(s^{(1)}, \dots, s^{(n)}\big), \qquad
    s^{(i)} = \big(\mathrm{digitize}(o_1, B_1^{(i)}), \dots, \mathrm{digitize}(o_d, B_d^{(i)})\big),
\]

and $\hat{Q}$ is stored as an $n \times b_1 \times \cdots \times b_d \times
|\mathcal{A}|$ table — one full aggregation table per tiling.

\subsection{Action selection and updates}

The action-value at $s$ is the mean over tilings,

\[
    \hat{Q}(s, a) = \frac{1}{n}\sum_{i=1}^n \hat{Q}^{(i)}\big(s^{(i)}, a\big),
\]

used directly for $\varepsilon$-greedy action selection. Each of the $n$
per-tiling tables is then updated independently with the same one-step
Sarsa error, so that the tiling that happens to distinguish two nearby
states most finely can still adapt on its own:

\[
    \hat{Q}^{(i)}\big(S_t^{(i)}, A_t\big) \leftarrow
    \hat{Q}^{(i)}\big(S_t^{(i)}, A_t\big) + \alpha \Big( R_{t+1} + \gamma \hat{Q}(S_{t+1}, A_{t+1}) - \hat{Q}^{(i)}\big(S_t^{(i)}, A_t\big) \Big), \quad i = 1, \dots, n.
\]

\begin{algorithm}[h]
\caption{Sarsa with tile coding}
\begin{algorithmic}[1]
\State Choose $n$ tilings and base bins $\mathbf{b}$; build $n$ jittered/offset aggregations
\State Initialize $\hat{Q}^{(i)}(s,a) \leftarrow 0$ for each tiling $i = 1,\dots,n$
\For{episode $= 1, 2, \dots$}
    \State $S_0 \leftarrow \big(S_0^{(1)}, \dots, S_0^{(n)}\big) \leftarrow$ tile-coded reset observation
    \State Choose $A_0$ $\varepsilon$-greedily from $\hat{Q}(S_0, \cdot) = \frac{1}{n}\sum_i \hat{Q}^{(i)}(S_0^{(i)}, \cdot)$
    \For{$t = 0, 1, 2, \dots$ until done}
        \State Take $A_t$, observe $R_{t+1}$ and tile-coded $S_{t+1}$
        \State Choose $A_{t+1}$ $\varepsilon$-greedily from $\hat{Q}(S_{t+1}, \cdot)$
        \For{$i = 1, \dots, n$}
            \State $\hat{Q}^{(i)}(S_t^{(i)}, A_t) \leftarrow \hat{Q}^{(i)}(S_t^{(i)}, A_t) + \alpha \big( R_{t+1} + \gamma \hat{Q}(S_{t+1}, A_{t+1}) - \hat{Q}^{(i)}(S_t^{(i)}, A_t) \big)$
        \EndFor
    \EndFor
\EndFor
\end{algorithmic}
\end{algorithm}

\section{State aggregation vs.\ tile coding}

\begin{center}
\begin{tabular}{@{}p{4.1cm}p{4.6cm}p{4.6cm}@{}}
\toprule
 & State aggregation & Tile coding \\
\midrule
Tables stored & $1$ & $n$ (one per tiling) \\
Effective resolution & bin width & finer than any single tiling's bin width \\
Generalization across nearby states & none (hard cell boundaries) & partial (shared tilings) \\
Memory & $\prod_i b_i \cdot |\mathcal{A}|$ & $n \cdot \prod_i b_i \cdot |\mathcal{A}|$ \\
Discontinuities at bin edges & all coincide, always visible & smoothed out (edges of different tilings rarely coincide) \\
\bottomrule
\end{tabular}
\end{center}

Both are still fundamentally tabular under the hood — every table entry is
updated independently by the same sample-average TD rule as ordinary
Sarsa — so all of the tabular convergence guarantees still apply
\emph{to the discretized problem}. What is lost is any guarantee about the
\emph{original} continuous-state problem: a coarser discretization can
introduce approximation error (aliasing states that should be
distinguished) that no amount of additional training removes, and picking
$\mathbf{b}$ and $n$ is itself a bias--variance-style tradeoff between
resolution/memory and how much data is needed to fill every cell.

\subsection{Beyond binning}

State aggregation and tile coding are the simplest members of a much
larger family of \textbf{function approximation} methods for continuous
(or just very large) state spaces, where $\hat{Q}(s,a;\theta)$ is
represented parametrically (linear-in-features, or a neural network as in
DQN) instead of as an explicit table, and the TD update becomes a gradient
step on $\theta$ rather than a per-entry table update. Tile coding is in
fact exactly linear function approximation with a particular fixed,
binary feature map $\mathbf{x}(s) \in \{0,1\}^{n \cdot \prod_i b_i}$ (one
active feature per tiling); it is presented here in tabular form because
the per-tiling tables can be updated independently with the ordinary
tabular Sarsa rule, without needing to introduce gradients explicitly.

\section{Reference}

Sutton, R.\ S., \& Barto, A.\ G.\ \emph{Reinforcement Learning: An
Introduction}, 2nd ed., Section~9.5.4 (Tile Coding).

\end{document}
